\documentclass[11pt,a4paper]{article}

% ---- Márgenes
\usepackage[a4paper, left=2cm, right=2cm, top=2.5cm, bottom=2.5cm]{geometry}

% ---- Idioma y fuentes (pdfLaTeX)
\usepackage[utf8]{inputenc}
\usepackage[T1]{fontenc}
\usepackage[spanish, es-tabla]{babel}

% ---- “Arial-like” con pdfLaTeX: Helvetica
\usepackage[scaled=0.92]{helvet}
\renewcommand{\familydefault}{\sfdefault}

% ---- Matemáticas y tablas
\usepackage{amsmath,amsfonts,amssymb}
\usepackage{tabularx,booktabs}
\usepackage{pdfpages}

% ---- Gráficos
\usepackage{graphicx}
\usepackage{subcaption}
\usepackage{float}
\usepackage[table,xcdraw]{xcolor}
\definecolor{verdePastel}{RGB}{204, 235, 197}
\definecolor{verdeFresco}{RGB}{60,160,60}

% ---- Estilo y utilidades
\usepackage{fancyhdr}
\setlength{\headheight}{14pt}
\usepackage{csquotes}
\usepackage{parskip}
\usepackage{array,longtable}
\renewcommand{\arraystretch}{1.4}
\newcolumntype{L}[1]{>{\raggedright\arraybackslash}p{#1}}
\newcolumntype{C}[1]{>{\centering\arraybackslash}p{#1}}
\newcolumntype{B}[1]{>{\raggedright\arraybackslash}p{#1}}

% ---- URLs más “rompibles”
\usepackage{xurl}

% ---- Interlineado 1.5
\usepackage{setspace}

% ---- Bibliografía (biber)
\usepackage[
  backend=biber,
  style=numeric,
  sorting=none,
  maxbibnames=99
]{biblatex}
\addbibresource{referencias.bib}

% ---- Hipervínculos (cargar al final)
\usepackage{hyperref}
\hypersetup{
  pdfauthor={Dustin Randel De Lara García},
  pdftitle={PEC3 Grupo B - Cuantización de la energía según la Física Cuántica },
  pdfkeywords={Física, Química, Energía, UNED},
  colorlinks=true,
  linkcolor=verdeFresco,
  urlcolor=verdeFresco,
  citecolor=verdeFresco
}

% ---- “Tabla” en lugar de “Cuadro”
\addto\captionsspanish{\renewcommand{\tablename}{Tabla}}

% ---- Encabezados/pies
\pagestyle{fancy}
\rhead{UNED}
\lhead{PEC3 Grupo B - Cuantización de la energía según la Física Cuántica}
\cfoot{Dustin Randel de Lara García \;--\; \href{mailto:ddelara3@alumno.uned.es}{ddelara3@alumno.uned.es}}
\rfoot{\thepage}

\begin{document}
\onehalfspacing

\subsection*{Descriptor}
Física y Química, 2.º Bachillerato (Científico/Tecnológico); grupo heterogéneo; 2 sesiones de 55 min. Tema: \textit{cuantización de la energía} ($E=h\nu$) y \textit{efecto fotoeléctrico}.
\subsection*{Objetivos de aprendizaje (al finalizar la sesión, el alumnado será capaz de):}
\begin{itemize}
  \item Enunciar e interpretar $E=h\nu$ (significado físico y unidades).
  \item Explicar cualitativamente el efecto fotoeléctrico: umbral $\nu_0$, frecuencia vs.\ intensidad.
  \item Resolver un problema tipo usando $K_{\max}=h\nu-\Phi$ y $eV_0=K_{\max}$.
\end{itemize}

\subsection*{Competencias y criterios (RD 243/2022, Anexo II, Física):}
Se trabajan principalmente las competencias específicas \textbf{CE1}, \textbf{CE3}, \textbf{CE5} y \textbf{CE6},
mediante los criterios de evaluación \textbf{1.2}, \textbf{3.1 - 3.3}, \textbf{5.1 - 5.2} y \textbf{6.1}
(asociados a resolución de problemas, uso riguroso de unidades/expresión científica, indagación con datos y
contextualización histórica de la física) \cite[págs.~143 - 144]{BOE_RD243_2022}. 

\subsection*{Contenido}
\textbf{1) Activación / motivación (5 min).} Pregunta: \emph{Sobre una placa de zinc, ¿por qué una luz roja puede no arrancar electrones, mientras que una luz violeta sí?} (recuperar idea de energía y frecuencia).

\textbf{2) Idea clave: cuantización (10 min).} Planck (1900): la energía no se intercambia de forma continua, sino en paquetes proporcionales a la frecuencia. Relación: $E = h\nu$, donde $h$ es la constante de Planck \cite{NIST_MeetConstants}. Breve contexto histórico (Planck 1900) \cite{MPG_PlanckPackets}.

\textbf{3) Einstein y el efecto fotoeléctrico (15 min).} Modelo de fotones aplicado (Einstein, 1905): cada electrón absorbe \emph{un fotón o ninguno}. Balance energético: $K_{\max}=h\nu-\Phi$.
Consecuencias observables:
\begin{itemize}
  \item Existe una \textbf{frecuencia umbral} $\nu_0=\Phi/h$ (si $\nu<\nu_0$, no hay emisión).
  \item La energía cinética máxima $K_{\max}$ depende de \textbf{frecuencia} $\nu$, no de la \textbf{intensidad} del haz de luz.
  \item No hay retardo apreciable (emisión prácticamente instantánea).
\end{itemize}
(Referencia histórica: Nobel a Einstein por el efecto fotoeléctrico) \cite{Nobel_Einstein1921}. Para una formulación clara y didáctica del fenómeno, consultar las siguientes fuentes: \cite{OpenStax_Photoelectric,Britannica_Photoelectric}.

\textbf{4) Actividad guiada con simulación en el navegador (15 min).} Simulador online: PhET \cite{PhET_Photoelectric}:
\begin{itemize}
  \item Fijar un metal (elige uno con $\Phi$ conocido en la propia simulación).
  \item Variar \textbf{frecuencia} y medir el \textbf{potencial de frenado} $V_0$.
  \item Tomar 4 - 6 puntos $(\nu, V_0)$ y representar mentalmente/rápidamente la tendencia lineal.
  \item Interpretación: $eV_0=h\nu-\Phi \Rightarrow V_0=\frac{h}{e}\nu-\frac{\Phi}{e}$, es una recta con pendiente $h/e$.
\end{itemize}

\textbf{5) Cierre (10 min).} Resumen de ideas:
\begin{itemize}
  \item \textbf{Frecuencia} controla \textbf{energía por fotón}. Intensidad controla \textbf{número de fotones}.
  \item Umbral $\nu_0$ y función de trabajo $\Phi$ como “peaje” energético.
  \item Enlace con dualidad onda-corpúsculo y física moderna (bloque curricular).
\end{itemize}

\subsection*{Duración}
\textbf{1 sesión de 55 minutos (Teoría)}:
5 min activación + 10 min cuantización + 15 min efecto fotoeléctrico + 15 min PhET + 10 min cierre.

\textbf{1 sesión de 55 minutos (Práctica/Evaluación)}:
15 min custionario rápido + 20 min mini-problema + 20 min PhET (recogida de datos e interpretación).
\subsection*{Evaluación}
\textbf{Instrumentos:}
\begin{itemize}
  \item \textbf{Cuestionario rápido} (5 ítems): 3 tipo test + 2 cortas (unidades de $h$, diferencia frecuencia/intensidad, definición de $\nu_0$).
  \item \textbf{Mini-problema} (1 ejercicio): calcular $K_{\max}$, $V_0$ o $\nu_0$ a partir de $\Phi$ y $\nu$.
  \item \textbf{Evidencia de indagación} (PhET): tabla de 4--6 medidas y 2 líneas de interpretación.
\end{itemize}

\textbf{Criterio de logro (simple y medible):}
\begin{itemize}
  \item Cuestionario: \textbf{$\ge 4/5$} correcto.
  \item Mini-problema: procedimiento correcto y resultado con unidades (se admite error numérico menor).
  \item PhET: datos coherentes + identifica que la tendencia lineal.
\end{itemize}

\textbf{Alineamiento CE--criterios--instrumentos (síntesis):}
\begin{center}
\small
\rowcolors{2}{verdePastel!35}{white}
\begin{tabularx}{\linewidth}{L{5.2cm}X}
\toprule
\textbf{CE y criterio(s) (RD 243/2022)} & \textbf{Instrumento + evidencia (criterio de logro)} \\
\midrule
CE3 (3.1) + CE6 (6.1) &
\textbf{Cuestionario rápido} (3 test + 2 cortas): explica cuantización ($E=h\nu$), umbral $\nu_0$ y diferencia
\emph{frecuencia/intensidad}. \textbf{Logro:} $\ge 4/5$. \\

CE1 (1.2) + CE3 (3.2--3.3) &
\textbf{Mini-problema}: aplica $K_{\max}=h\nu-\Phi$ y $eV_0=K_{\max}$, con \textbf{unidades} y pasos razonados.
\textbf{Logro:} procedimiento correcto + unidades (error numérico menor admisible). \\

CE5 (5.1--5.2) + CE4 (4.1) &
\textbf{PhET}: tabla 4--6 pares $(\nu,V_0)$ + 2 líneas de interpretación (“$V_0$ crece con $\nu$”, tendencia lineal).
\textbf{Logro:} datos coherentes + conclusión correcta. \\
\bottomrule
\end{tabularx}
\end{center}
\newpage
\section*{Bibliografía}
\printbibliography[heading=none]
\newpage
\includepdf[pages=1]{DECLARACÓN_JURADA_DE_AUTORÍA.pdf}
\end{document}
